\documentclass[11pt]{article}
\usepackage[margin=1in]{geometry}
\usepackage{setspace}
\usepackage{hyperref}
\hypersetup{colorlinks=true, allcolors=black}
\setlength{\parindent}{0pt}
\setlength{\parskip}{0.5em}
\onehalfspacing
\pagestyle{empty}

\begin{document}

\begin{center}
{\Large Teaching Statement}\\[0.3em]
HSU Ke Cheng
\end{center}

I have taught at the Hong Kong University of Science and Technology (HKUST) in several roles. From 2020 to 2025, I was a teaching assistant for undergraduate macroeconomics, the economics of uncertainty and information, and PhD microeconomic theory, leading tutorials, holding office hours, and grading. Through office hours, I learned to identify where a student's reasoning becomes unclear and rebuild the argument from there. In the doctoral course, I often approached the same question in several ways until the idea became clear. My average evaluation across these five appointments was 4.72 out of 5. I also lectured in the math camp for the MSc in Economics in 2026, designed the lecture material myself, and built the mathematical background step by step for students from a non-math background, receiving 4.53 out of 5. As Teaching Assistant Coordinator for the Business School from 2025 to 2026, I helped postgraduate students design teaching materials and evaluated their mock teaching.

I am prepared to teach microeconomic theory and industrial organization from the undergraduate to the PhD level, and information economics and mathematics for economists through the master's level. At the undergraduate level, I can also teach macroeconomics and econometrics. Beyond these theoretical courses, I am capable of teaching an introductory course in computational methods. I have substantial experience in structural estimation using R, Python, and C++, with Git for version control and collaboration. In my research, the estimation code was developed in R, while Python supported the data collection and OCR workflows that digitize scanned and historical sources. I also use AI to improve my efficiency in coding, debugging, and day-to-day workflow. This background prepares me to teach a course on conducting research with AI and managing workflows efficiently.

As the instructor of my own courses, I would keep the same emphasis on intuition, active reasoning, and clear explanation. I want students to use economic models not simply to solve equations, but to distinguish an observed pattern from the mechanism that produces it. A new restaurant opening on campus illustrates this: students may treat a long line as evidence that the food is good, but the line is only a signal, and as more people observe it and join, the line can keep growing even when the food is bad. This familiar example gives students a way into herding under imperfect information. Another example asks why a government might issue consumption vouchers instead of cash: because vouchers must be spent, more of the transfer enters current demand and the multiplier is larger. These examples shape how I teach. I begin with the intuition and introduce the formula only after students understand the question it answers. Once the intuition is clear, I ask students to work through the reasoning and tell me where it stops being clear. When one student asks a question, it often helps several others and shows me where to explain more simply. Teaching is most fulfilling when a student who was stuck begins to see why an economic outcome occurs, and finding that clearer explanation is the part of teaching I enjoy most.

\end{document}
